\documentclass[11pt]{article}
\usepackage[margin=1in]{geometry}
\usepackage{amsmath}
\usepackage{booktabs}
\usepackage{array}
\usepackage{enumitem}
\usepackage{titlesec}
\usepackage{parskip}
\usepackage[hidelinks]{hyperref}
\usepackage{microtype}

\titlespacing*{\section}{0pt}{8pt}{4pt}
\titlespacing*{\subsection}{0pt}{6pt}{2pt}
\titleformat{\section}{\normalfont\large\bfseries}{\thesection}{0.6em}{}
\titleformat{\subsection}{\normalfont\bfseries}{\thesubsection}{0.6em}{}

\setlist{itemsep=1pt, topsep=2pt, parsep=0pt}
\setlength{\parskip}{2pt}
\setlength{\tabcolsep}{5pt}
\renewcommand{\arraystretch}{0.95}

\title{\vspace{-1em}Design Note --- Learning from Click-Logs on EB-NeRD and MIND\vspace{-0.5em}}
\author{\small CS4.406 Information Retrieval \& Extraction --- Assignment 2}
\date{\small \today}

\begin{document}
\maketitle
\vspace{-1.5em}

\section{What We Built}

Assignment 1 built lexical (BM25) and semantic (embedding) candidate
generation for both datasets. Assignment 2 adds the behavioural layer: engineered
click-history features, a two-stage retrieve-then-rank pipeline, and --- the
graded gap A1 left open --- a literal reproduction of each dataset's official
NRMS baseline, improved by one principled change and tested with a paired
bootstrap.

\subsection*{1.1 Click-history and session features (Q1)}

EB-NeRD already had 65 engineered features (Candidate K, ADR-013): long/short-term
category/topic/embedding affinity, freshness (\texttt{article\_age\_h}, the single
top feature by gain), and --- added this assignment for Q1 --- session features
(\texttt{session\_position}, \texttt{session\_start\_gap\_h}, three
\texttt{context\_*} match features). MIND had none of this: only frozen-embedding
similarity and symbolic category/entity overlap (Candidates A--J).

Candidate L (A2 Q1, this assignment) adds five engineered features for MIND: click
count, category match, subcategory match (the last two already existed unweighted
in \texttt{features.py}), and two new \textbf{recency-weighted} affinity features
using the same geometric decay (decay$=0.9$, most-recent-last) already established
for MIND's embedding-query recency experiment (Candidate C). Scored by a cheap
CPU-only logistic-regression combiner --- deliberately not a new NRMS run, so it
does not contend with the treatment training job below.

\textbf{Result: a real, CI-clear loss.} AUC $0.6176$ [$0.6153,0.6197$] vs.\ the
deployed embedding baseline's $0.6340$ --- paired $-0.0164$ [$-0.0191,-0.0139$],
and every metric (MRR, nDCG@5, nDCG@10) loses the same way (ADR-016). Reported
plainly, not explained away: no hyperparameter search was run, and the CI
half-width ($\sim$0.001--0.002) is far smaller than the 0.016 gap, so tuning
would not close it. Click-history statistics alone --- recency-weighted or not
--- do not carry enough per-candidate signal to beat frozen embedding
similarity on MIND, consistent with Candidate D's earlier standalone loss
(0.613) and with Candidate F's finding that even a 7-feature combiner
including embeddings only barely avoided losing overall. Not adopted; MIND's
real win remains Candidate J and, per this assignment, the reproduced-then-
improved official NRMS (\S2.2).

\textbf{Stated as a dataset limitation, not a gap to force-fill:} MIND has
\emph{no} freshness feature and \emph{no} session feature. Confirmed by direct
measurement against the processed data, not assumed --- \texttt{article\_published\_time}
is null for all 72{,}023 MINDlarge articles, and \texttt{session\_id}\slash\texttt{dwell\_time}\slash
\texttt{scroll\_percentage}\slash\texttt{is\_front\_page} are null for all 14{,}085{,}557
impression rows. \texttt{src/datasets/mind.py} hardcodes these columns to \texttt{None}
at parse time because MIND's source TSVs never carry this information at all --- it
is not a sampling artifact and cannot be engineered around.

\textbf{A modeling fact worth stating plainly:} click count is identical for every
candidate within one impression (it depends only on the user). A linear combiner's
contribution to a per-impression ranking metric (AUC/MRR/nDCG@K --- all rank-order-only)
is provably invariant to an additive per-impression constant. Rather than asserting
this from theory alone, the experiment script includes an ablation that zeroes the
feature and re-scores every dev impression: the maximum absolute AUC change was
\textbf{exactly $0.00\mathrm{e}{+00}$} --- not approximately zero, bit-for-bit zero,
matching the proof precisely. Click count is still reported, per Q1's explicit ask,
with this caveat attached rather than hidden; its fitted coefficient (0.064) is
nonzero because it can still reduce the pointwise training loss (which differs
across users) even though it cannot move any impression's internal ranking.

\textbf{Leakage.} MIND provides no per-click timestamps, so the EB-NeRD-style
temporal-boundary check cannot be built for MIND. The structurally equivalent guard
the dataset actually supports is reused instead: since a user's history is one fixed
snapshot for the whole split, every history-derived feature must trace back to
\emph{exactly one} value regardless of which of that user's real impressions produced
it. A new integration test exercises this against real MINDsmall data (not a
synthetic call-the-function-twice check, which would prove only determinism) and
confirms candidate-derived components legitimately vary while history-derived ones
do not.

\subsection*{1.2 Two-stage retrieve-then-rank, and an honest retrieval ceiling (Q2)}

Both leaderboard submissions rank each impression's own shown candidate list ---
what the competitions actually score. A2 additionally asks for top-$K$ retrieval
from A1's own generator, $K\sim100$--$200$. We report this honestly, separately:
A1's measured recall@200 is 2.62--2.78\% on MIND and 1.21--2.77\% on EB-NeRD, so
at least $\sim$97\% of clicked articles never enter a top-200 list from either
retriever. No re-ranker can recover a click the candidate generator never surfaced.

\textbf{Only NRMS (both datasets) re-ranks in that evaluation.} Candidate K's GBDT
is excluded by construction: its features are \emph{impression-conditional}
(\texttt{position\_in\_view}, the five \texttt{*\_rank\_in\_imp} features,
\texttt{n\_candidates}, three \texttt{context\_*} features), defined relative to the
other candidates the publisher actually showed. A retrieved-but-never-shown
candidate has no position in that list and no in-list rank to compare against ---
any value would be invented and would silently drive the ranking. A reduced
``retrieval-safe'' GBDT was considered and rejected: it would answer a different
model's ceiling, not Candidate K's, at the cost of a second model to maintain.

\section{Baseline Reproduced, Then Beaten (Q3) --- A/B Framing}

\textbf{Control} = the official baseline, reproduced. \textbf{Treatment} = control
plus exactly one principled change. \textbf{Test} = paired bootstrap 95\% CI of
(treatment $-$ control) over the identical impressions, resampling users --- the
same statistic ADR-010's Candidate G introduced. \textbf{Guardrails} that must not
regress even if the primary metric wins: diversity@10, novelty@10, coverage@10, and
serving latency (p99, \S3).

\textbf{Reproduction method.} The brief names the official code directly:
ebnerd-benchmark's NRMS for EB-NeRD, \texttt{microsoft/recommenders}' NRMS for MIND.
We attempted this (Option A) on Ada within a one-session timebox; it did not train
--- every reachable machine's PyPI throughput (0.09--0.37\,MB/s) made the $\sim$3\,GB
TensorFlow+CUDA install multi-hour, and one job held a GPU idle $\sim$4h before we
diagnosed it. Per the pre-agreed timebox rule, we fell back to \textbf{Option B}: a
layer-for-layer PyTorch port of the official Keras graph (bias-free Q/K/V attention
at the official 20$\times$20 width, no output projection, no padding mask, the exact
\texttt{AttLayer2} pooling formula, official hyperparameters read from the shipped
config files rather than retyped) plus the official loaders' exact input semantics
(most-recent-50 history, MIND's zero-pad negative sampling, EB-NeRD's
with-replacement sampling). This must be read as \emph{a faithful reimplementation},
never as ``the official code'' --- and it is reported as such throughout, with the
revisit condition (a cheap TF environment) recorded rather than silently dropped.

\subsection*{2.1 EB-NeRD --- complete}

\begin{center}
{\footnotesize
\begin{tabular}{@{}lccccccc@{}}
\toprule
& AUC & 95\% CI & MRR & nDCG@5 & nDCG@10 & Div.@10 & Nov.@10 \\
\midrule
Control (reproduced NRMS) & 0.5613 & 0.5597--0.5629 & 0.3512 & 0.3904 & 0.4680 & 0.7894 & 17.208 \\
Treatment (+freshness) & \textbf{0.5677} & 0.5659--0.5696 & 0.3555 & 0.3957 & 0.4725 & 0.7900 & 17.236 \\
\bottomrule
\end{tabular}\par}
\end{center}

\textbf{Paired treatment $-$ control:} AUC $+0.0064$ [$+0.0053$, $+0.0075$]; MRR
$+0.0043$; nDCG@5 $+0.0053$; nDCG@10 $+0.0044$ --- all CI-clear. \textbf{No guardrail
regressed}: diversity $+0.0006$ and novelty $+0.0280$, both CI-clear improvements;
coverage fell $0.2021\to0.1987$ ($-1.7\%$ relative), reported as a point difference
only (ADR-007: a set-union statistic is mechanically biased under bootstrap
resampling), flagged rather than hidden.

Two things the headline must not obscure. \textbf{The gain is small and merely
\emph{tightly} estimated} --- the narrow CI reflects $n=244{,}647$ impressions, not
a large effect. \textbf{The improved baseline (0.5677) is still far below Candidate
K (0.7588)} on the identical split; A1's GBDT remains the strongest EB-NeRD system
this project has, and beating the official NRMS baseline does not change that.

\textbf{Head/tail slicing (Q5) reverses the headline.} An impression's clicked
article is \emph{head} if it sits in the top 20\% of train-split click counts,
else \emph{tail} (train-only, so no dev-click leakage; never-clicked-in-train
articles fall in tail by construction). The paired gain is entirely a tail effect:

\begin{center}
{\footnotesize
\begin{tabular}{@{}lrrl@{}}
\toprule
Slice & $n$ & Paired $\Delta$AUC (95\% CI) & Verdict \\
\midrule
Head (top 20\% by train clicks) & 9{,}541 & $-0.0217$ [$-0.0258$, $-0.0178$] & CI-clear \textbf{loss} \\
Tail & 235{,}106 & $+0.0076$ [$+0.0064$, $+0.0087$] & CI-clear gain \\
\bottomrule
\end{tabular}\par}
\end{center}

The $+0.0064$ aggregate is a tail gain diluted by a real head regression on 3.9\%
of impressions --- exactly the failure mode slicing exists to catch. Mechanistically
plausible: a freshness-weighted ranker should help least where an article is
popular enough to be clicked regardless of age.

\textbf{Anti-gaming (Q9): pricing the serving-time-unavailable features.}
\texttt{context\_read\_time}/\allowbreak\texttt{context\_scroll\_percentage} were flagged in
ADR-013 as EB-NeRD's weakest serving-time claim (a page's total read time is only
known once the user leaves it). A fourth GBDT arm withholding both gives
$0.7570$ vs.\ $0.7581$ with them --- paired $+0.0011$ [$+0.0008,+0.0015]$, CI-clear
but negligible against the $+0.2140$ margin over \texttt{embed\_sim}. For contrast,
withholding \emph{position} features \emph{improves} AUC by $0.0007$. Neither
flagged group is load-bearing; freshness is (\texttt{article\_age\_h} remains top
by gain). \texttt{K\_rank} and \texttt{K\_rank\_nopos} reproduced bit-identically
across the three- and four-arm runs, incidental evidence of a deterministic pipeline.

\subsection*{2.2 MIND --- complete}

\begin{center}
{\footnotesize
\begin{tabular}{@{}lccccc@{}}
\toprule
& AUC & 95\% CI & MRR & nDCG@5 & nDCG@10 \\
\midrule
Control (reproduced NRMS, 10 ep.) & 0.6831 & 0.6821--0.6839 & 0.3803 & 0.3620 & 0.4279 \\
Candidate J (A1's best, same split) & 0.6579 & --- & 0.3581 & 0.3411 & 0.4064 \\
Published NRMS (paper, \emph{test} split) & 0.6776 & --- & 0.3305$^\dagger$ & 0.3594 & 0.4163 \\
\bottomrule
\end{tabular}\par}
\end{center}
$^\dagger$ Not directly comparable: the official evaluator sums $1/\text{rank}$ over
\emph{all} clicked items; ours is first-hit-only (established in A1's Codabench
converter validation).

The reproduced control \textbf{beats Candidate J by $+0.0252$ AUC} on the identical
MINDlarge-dev split (marginal, not paired --- J's per-impression scores were never
saved) and lands close to the published test-split figure, which is the evidence
Q3 needs that the reproduction is trustworthy. Head/tail and warm/cold slices are
already computed for the control: head 0.7411 vs.\ tail 0.6474; warm 0.6932 vs.\
cold 0.6209 --- a 0.0723 gap, marginally \emph{wider} than A1's embedding-baseline
gap (0.0682), so the reproduced baseline lifts both cohorts $\sim$0.05 without
closing cold-start, echoing ADR-008's conclusion for embeddings.

\textbf{Treatment: title+abstract input} (title 30 tokens + abstract 50, widths
from the official \texttt{nrms.yaml}/\texttt{naml.yaml}), ablated against the
title-only control. \emph{The original treatment spec --- GloVe vs.\ random
init --- was reconsidered before running it}: the official NRMS is already
GloVe-initialised, so that comparison would have measured the control against a
deliberately degraded control, not an improvement. Training took 34.06h (10
epochs, no errors), comfortably inside the 72h wall limit.

\begin{center}
{\footnotesize
\begin{tabular}{@{}lccccccc@{}}
\toprule
& AUC & 95\% CI & MRR & nDCG@5 & nDCG@10 & Div.@10 & Nov.@10 \\
\midrule
Control (reproduced NRMS) & 0.6831 & 0.6821--0.6839 & 0.3803 & 0.3620 & 0.4279 & 0.8343 & 17.623 \\
Treatment (+title\slash abstract) & \textbf{0.6868} & 0.6859--0.6877 & 0.3829 & 0.3691 & 0.4314 & 0.8283 & 17.643 \\
\bottomrule
\end{tabular}\par}
\end{center}

\textbf{Paired treatment $-$ control:} AUC $+0.0037$ [$+0.0031,+0.0044$]; MRR
$+0.0026$; nDCG@5 $+0.0071$; nDCG@10 $+0.0034$ --- all CI-clear wins.
\textbf{A guardrail genuinely regressed:} diversity $-0.0060$
[$-0.0062,-0.0057]$, CI-clear, while novelty improved $+0.0197$ (CI-clear) and
coverage was flat ($-0.0001$, point only). \textbf{This is the first guardrail
regression in this project's entire history} --- every earlier win, including
EB-NeRD's freshness treatment above, improved every CI-bearing guardrail. Per
the professor's explicit framing, a guardrail regression is not something a
primary-metric win is allowed to paper over: it is reported as a real cost of
the change, not smoothed into the accuracy headline. Plausible mechanism:
49 additional abstract tokens likely sharpen the model's confidence toward
specific topical matches at some expense to the categorical variety of what it
ranks highly.

\textbf{Head/tail slicing independently reproduces EB-NeRD's exact directional
pattern} --- a different dataset, a different treatment, a different mechanism,
the same qualitative result: head $-0.0063$ [$-0.0071,-0.0053$] (CI-clear loss,
$n=143{,}164$) vs.\ tail $+0.0098$ [$+0.0090,+0.0108$] (CI-clear gain,
$n=233{,}307$). Both treatments help most where the content signal is weakest
relative to prior popularity and can mildly hurt where an article is already
identifiable from its title alone --- independent replication, not designed in
advance. \textbf{The gain is entirely a warm-user effect}: warm $+0.0043$
[$+0.0036,+0.0050]$ CI-clear, cold $+0.0002$ [$-0.0020,+0.0022]$ \emph{not
significant} --- sensible, since richer candidate text cannot improve a user
representation built from zero history. The reproduced-then-improved pipeline
still does not close cold-start, consistent with every earlier finding in this
project.

\section{Serving and Scale Analysis (Q4)}

Everything below is a measured \emph{single-request} p99, not the batched,
amortised figure this project's stage-level profiling (ADR-014) originally
reported --- the gap between the two is itself a finding (below).

\begin{center}
{\footnotesize
\begin{tabular}{@{}lrrrr@{}}
\toprule
Instance & Index/store & p99 & SLA 100ms & \$/1k (1 worker) \\
\midrule
MIND / CPU (c6i.2xlarge) & 164.0\,MB & \textbf{94.67\,ms} & PASS, $1.1\times$ & \$0.0031 \\
MIND / GPU (g4dn.xlarge) & 164.0\,MB & \textbf{33.15\,ms} & PASS, $3.0\times$ & \$0.0007 \\
EB-NeRD (\texttt{nopos}) / CPU & 190.5\,MB & \textbf{22.74\,ms} & PASS, $4.4\times$ & \$0.0009 \\
\bottomrule
\end{tabular}\par}
\end{center}

Both MIND rows are genuinely measured on their own hardware, not one profile
priced onto both instance types --- an early version of the cost script did
exactly that (excluded rather than reported the invalid row), and a second,
mirrored version of the same bug (this time on the CPU row, once the GPU
profile became the newest file) was caught and fixed with a regression test
before these numbers were trusted.

\textbf{Batched vs.\ single-request differs by mechanism, and by how much room
there is to amortise.} EB-NeRD's $\sim$30$\times$ gap (0.33\,ms/impression
amortised vs.\ $\sim$10\,ms for one cold request) comes from per-call
vectorisation overhead paid once per 5{,}000-impression chunk instead of once
per impression. Quoting the batched figure as request latency would have
understated EB-NeRD's real serving cost by a factor of thirty. MIND's batched
sweep instead hoists per-\emph{user} work (tokenize, BM25/embedding scoring)
out of the impression loop --- but the real MIND population averages only
1.47 impressions/user (ADR-014), so there is little to amortise, and the
batched figure (40.46\,ms/impression) and the cold single-request mean
(33.29\,ms) sit within $\sim$1.2$\times$ of each other, in the \emph{opposite}
direction from EB-NeRD's gap. Reported as measured, not smoothed into a false
parallel between the two datasets.

\textbf{On CPU, MIND's p99 headroom is thin, and its observed max (126\,ms)
already breaches the 100\,ms target} --- NRMS is 78\% of a request's cost there.
\textbf{On GPU, the same target has real headroom} ($3.0\times$): NRMS drops to
36\% of the request, matching this project's own GPU-vs-CPU finding elsewhere
in ADR-014 (NRMS $8.98\times$ faster on a 2080\,Ti). EB-NeRD is comfortable on
CPU alone (4.4$\times$) and dominated by feature-frame construction, not the
booster, so a GPU figure was not pursued for it.

The \textbf{1-worker} \$/1k figures above are the conservative end of a range:
an \emph{optimistic} column (not shown) assuming perfect linear scaling across
an instance's vCPUs would give roughly 8$\times$ lower figures, but this
project's own profiling contradicts that assumption directly (a measured
$2.5\times$ slowdown from a single concurrent job on one 8GB machine,
ADR-014) --- real per-instance throughput is sublinear, so the 1-worker number
is the more honest one to lead with. Prices are AWS on-demand list, us-east-1,
captured 2026-09-12.

\section{Where the System Breaks at 10$\times$}

\textbf{MIND breaks on CPU latency, not memory --- and GPU serving is a
measured fix, not a hoped-for one.} The index (164\,MB) scales with
articles/vocabulary, not users --- MINDlarge-test's 120,959 articles project to
$\sim$250\,MB, still trivial. On CPU, p99 headroom is only $1.1\times$ with
NRMS at 78\% of a request; \textbf{on GPU, the identical request-level
benchmark measures $3.0\times$ headroom} (\S3), not a projection from a
training-time speedup. A second, complementary fix needs no new hardware at
all: this project's own NRMS candidate-vector cache measures $7.02\times$ on
CPU with no accuracy cost (ADR-014). At 10$\times$ load, GPU serving buys real
room but not unlimited room --- $3.0\times$ headroom absorbs roughly a
3$\times$ traffic increase before the SLA is threatened again, so 10$\times$
load still needs horizontal scaling (more GPU workers) on top of the device
change, not instead of it.

\textbf{EB-NeRD breaks on memory, specifically user profiles.} They cost
5.35\,KB/user (82.1\,MB for 15{,}342 users); the article table (108.2\,MB)
scales with articles instead. At \texttt{ebnerd\_large}'s real 791{,}582 users,
the same per-user rate projects to $\sim$4.2\,GB of profiles alone --- from a
190\,MB feature store today to something that no longer sits comfortably beside
a model on a small serving instance. Latency is not the constraint here
($4.4\times$ headroom).

\textbf{The treatment itself has a real 10$\times$-adjacent cost story.} MIND's
title+abstract treatment costs $2.98\times$ the control per training step ---
measured, not assumed, from the first 45{,}000 of 105{,}740 steps before
committing the full $\sim$37h job. The same accounting applies to any future
feature addition: verify the marginal training cost on a small slice before
paying for the whole run, which is exactly how the 37h estimate was validated
here rather than discovered as an overrun.

\section{Limitations and Open Items}

\begin{itemize}
\item Q9's serving-time-unavailable ablation and Candidate K's 65-feature
  retrain are measured on \texttt{ebnerd\_small}, not the \texttt{ebnerd\_large}
  split behind the real 0.7542 Codabench score --- that bundle is unreachable
  at Ada's measured $\sim$18\,KB/s S3 rate with the cluster home directory's
  remaining $\sim$4.9\,GB.
\item Reproducibility is bitwise on CPU (verified: the reproduced-NRMS smoke
  tests and the EB-NeRD control's independent re-run matched to every printed
  digit) but
  only stable to $\sim$8 decimals run-to-run on GPU (ordinary CUDA reduction
  non-determinism, confirmed not to be a code defect) --- stated because this
  project's principle is ``same code + same data = same results,'' and the
  GPU qualification is real, not glossed over.
\end{itemize}

\end{document}
